% =============================================================================
%  md2pdf —— cn 档位（ctexart）的页眉修正
%
%  为什么需要这个文件：
%  ctexart 的 headings 页眉会把章标题整条交给 \MakeUppercase。对纯中文标题无所谓，
%  但技术文档的标题里常有拉丁词（Markdown、API、PDF），大写之后就成了 "MARKDOWN"，
%  同一份文档在正文里写 Markdown、在页眉里变 MARKDOWN，看起来像是排版出错。
%
%  press 档位早就用 \nouppercase{\rightmark} 处理过这件事（见 press.tex）。
%  这里按同一套模式给 cn 档补上，保持两档行为一致。
% =============================================================================

\usepackage{fancyhdr}
\setlength{\headheight}{15pt}
\pagestyle{fancy}
\fancyhf{}
% 左页眉：章标题，原样保留大小写；右页眉：页码
\fancyhead[L]{\small\nouppercase{\rightmark}}
\fancyhead[R]{\small\thepage}
\renewcommand{\headrulewidth}{0.4pt}
\renewcommand{\footrulewidth}{0pt}
% 页眉只放标题文本，不放 \thesection：
% pandoc 在 --no-numbers 时用 \setcounter{secnumdepth}{-\maxdimen} 关掉编号，
% 此时 \thesection 仍是 0，写进页眉会得到 "0.<n>" 这种没有意义的编号。
% 标题本身该带编号的（如本手册的 "4. 写作：…"）写在标题里就行。
\renewcommand{\sectionmark}[1]{\markright{#1}}
\renewcommand{\subsectionmark}[1]{}
